%!TEX root = ../thesis-trier.tex



Have you ever wanted to start all over again? Has your smartphone or computer ever frozen and you had to restart it? 
If so, then you are familiar with the idea of resetting things,
the central concept of the present work.  
More specifically, it is about resetting (also called synchronizing) 
finite state transition systems, also called finite automata. 
 Finite automata process input sequences (also called words) and can be used to model a wide variety of systems, for example,
switching circuits, networks, or simple
behavior patterns in artificial intelligence. 
In all these situations, knowing whether we can reset the system is of interest. A circuit component might be out of sync with its environment, can we reset it? 
A robot might be stuck, or we are simply unsure about its current state,
so can we reset its state of operation? 
However, while this problem is tractable
for completely specified and deterministic automata,
only a small deviation from this ``optimal'' situation, i.e., having only partial information, allowing nondeterminism, or constraining the admissible input
sequences, results in intractable
algorithmic problems.

Here, we are concerned with one particular problem variant: 
constraining the admissible input
sequences (the \emph{constrained synchronization problem}).
In this setting, what computational complexity can we deduce 
from the kind of constraint put on the input sequences and/or
types of input automata?
We state various results on how to manipulate sets of admissible synchronizing input sequences
so as to retain a certain complexity of the problem.
We also identify various conditions 
that either yield tractable or intractable problems.
More specifically, the following complexity classes will arise: 
$\AC^0$, \LOGSPACE, \NL, \PTIME, \NP, \PSPACE.
Except for~$\AC^0$ and~\PTIME, we 
give sets of admissible synchronizing sequences
yielding constrained problems that are complete for these classes.
Moreover, we investigate the following classes
of admissible sets of synchronizing sequences: finite, infinite, sparse, letter-bounded, commutative, or
constraints specified by certain small-state automata; where we give
a complete complexity classifications for the latter three classes.
Lastly, we  consider the following types of input automata to the problem:
monotonic, partially ordered, commutative, solvable, SI- and $\rho$--automata.
We obtain, for example, that for commutative input automata, the problem is always in \PTIME.

Now, the set of synchronizing 
sequences, i.e., our search space, of a given input automaton
is acceptable by a finite automaton as well.
This leads naturally to the investigation of accepting automata for sets
of synchronizing sequences, specifically for the question of when they
become largest possible.
Here,  we identify and investigate two conditions
that control the size of a minimal automaton
for the set of synchronizing sequences:
reachability and distinguishability.
For both notions, we establish various structural results and connections with concepts
from the literature such as, for example, completely reachable automata.
Decision problems and algorithmic questions 
that naturally arise
from these notions are also investigated.
We state conditions yielding extremal minimal automata accepting sets of synchronizing sequences.
Moreover, we identify a connection to permutation group theory, which 
is linked to automata theory by the group of units
of transformation monoids of automata.
Specifically, via graph-theoretical constructions,
we obtain new characterizations of primitive permutation groups
based on associated synchronizing automata.
Our results link permutation groups and vertex-transitive graphs
to the theory of synchronizing automata.
Furthermore, we explore relations of the size of automata accepting sets of synchronizing
sequences to the length of shortest synchronizing sequences.